\begin{opinion}
   
El fraude es un problema recurrente en la educación. Los profesores luchan constantemente contra las faltas morales tanto en exámenes como en proyectos y evaluaciones teóricas. Nuestra facultad no está exenta de este problema. La asignatura de Programación, de primer año de Ciencias de la Computación en MATCOM, tiene un sistema en el que el 50\% de la calificación del estudiante depende de proyectos evaluativos. Es común que año tras años los profesores encontremos, a ojo, proyectos clonados que representan faltas éticas. Sobre este marco, la estudiante desarrolló un software que nos resultará altamente útil.

Un \textit{software} capaz de detectar plagios en proyectos de C\# tendrá dos claras aportaciones a nuestra carrera. En primer lugar, evitará que prácticas inmorales pasen desapercibidas. En segundo lugar, automatizará una tarea que consume valioso tiempo a nuestro claustro: Buscar estas prácticas amorales.

María hizo una amplia investigación del estado del arte actual y notó que las herramientas libres de costo para comparación de código en lenguaje C\# (lenguaje que se usa en nuestra asignatura) eran ineficientes, casi inexistentes. Ella fue capaz de editar una gramática de C\# 6.0 para convertirla en una 12.0. Luego utilizó esa gramática para extraer características de los AST de un amplio conjunto de códigos de entrenamiento. Finalmente utilizó estas características para entrenar un modelo clasificador de aprendizaje automático. Este trabajo demostró conocimientos en un gran abanico de campos técnicos como: Compilación, Inteligencia artificial y programación orientada a objetos.

Como tutores, estamos orgullosos de que María haya convertido una queja aleatoria de un profesor de programación: "Qué complicado estar buscando trampas por más de 100 proyectos", en un excelente trabajo investigativo.

Por todo lo anterior creemos que estamos en presencia de una excelente científica de la computación, que al igual que el trabajo, merece la máxima calificación.\\
Rodrigo García Gómez y Alberto Fernández Oliva
   
\end{opinion}